\chapter{Supervised motion and measured limits}
\label{ch:hwmotion}

This chapter reports what the robot actually did. Every number has a date and a
session behind it.

\section{What passed, and when}

\begin{table}[h]
\centering
\small
\begin{tabularx}{\linewidth}{@{}llX@{}}
\toprule
\textbf{Gate} & \textbf{Date} & \textbf{Result} \\
\midrule
Bridge, non-motion & 2026-07-22 & Bridged topics at steady rates; robot state readable. \\
Action shims, interlocks & 2026-07-22 & Every malformed and unsafe goal rejected, zero motion. \\
Supervised motion & 2026-07-24 & Both arms, out and back, worst tracking error 0.0078\,rad. \\
MoveIt execution & 2026-07-24 & Planner success, 31 waypoints, settled 0.0128\,rad from target. \\
Speed characterisation & 2026-07-25 & Path-tolerance ceiling found; see below. \\
MoveIt re-run & 2026-07-25 & Mostly aborted on the left arm; see Section~\ref{sec:moveit-hardware}. \\
\bottomrule
\end{tabularx}
\end{table}

All of it on one robot, supervised. None of it establishes sustained duty,
gripper commands, or behaviour on a second machine.

\section{The speed ceiling}
\label{sec:speed}

The first supervised motion ran at 0.117\,rad/s against a shim clamp of
2.0\,rad/s. That is a factor of seventeen of unused headroom, which means no
tolerance figure in the workspace said anything at all about speed. The
2026-07-25 session escalated deliberately, one step at a time, with the recorder
running.

\begin{table}[h]
\centering
\small
\begin{tabular}{@{}rrrrl@{}}
\toprule
\textbf{Step} & \textbf{offset (rad)} & \textbf{duration (s)} & \textbf{rad/s} & \textbf{peak lag} \\
\midrule
1 & 0.35 & 3.0  & 0.12 & 0.037\,rad \\
2 & 0.35 & 1.5  & 0.23 & 0.094\,rad \\
3 & 0.50 & 1.0  & 0.50 & \textbf{0.198\,rad $\rightarrow$ abort} \\
4 & 0.50 & 0.5  & 1.00 & not reached \\
5 & 0.50 & 0.35 & 1.43 & not reached \\
\bottomrule
\end{tabular}
\end{table}

It stopped at step 3 with \texttt{left\_s1 lags its setpoint by 0.202 rad (limit
0.200)}. The arm then held to 0.0004\,rad over six seconds --- the abort-and-hold
path worked exactly as designed.

\begin{figure}[h]
\centering
% Measured on BR-01, 2026-07-25 (I20). Three points, one straight line, and the
% reason path_tolerance_rad is a speed limit rather than a safety margin.
% Source: docs/hardware_day_plan.md P4a escalation table.
\begin{tikzpicture}
\begin{axis}[
  width=0.82\linewidth, height=6.2cm,
  xlabel={commanded velocity (rad/s)},
  ylabel={peak tracking lag (rad)},
  xmin=0, xmax=1.1, ymin=0, ymax=0.46,
  grid=major, grid style={coderule!60},
  legend pos=north west,
  legend style={font=\scriptsize, draw=coderule},
  label style={font=\small}, tick label style={font=\scriptsize},
]

% lag = 0.4 s * velocity -- the arm runs a constant time behind its setpoint.
\addplot[domain=0:1.1, samples=2, thick, noterule, dashed]
  {0.4*x};
\addlegendentry{$\mathrm{lag} \approx 0.4\,\mathrm{s} \times v$}

% The tolerance that stops the move.
\addplot[domain=0:1.1, samples=2, thick, warnrule]
  {0.2};
\addlegendentry{\texttt{path\_tolerance\_rad} $=0.2$}

\addplot[only marks, mark=*, mark size=2.2pt, black]
  coordinates {(0.12,0.037) (0.23,0.094)};
\addlegendentry{completed}

\addplot[only marks, mark=triangle*, mark size=3.4pt, warnrule]
  coordinates {(0.50,0.198)};
\addlegendentry{aborted, \texttt{left\_s1}}

\node[font=\scriptsize, align=left, anchor=west, text=comment]
  at (axis cs:0.56,0.28) {shim clamp 2.0\,rad/s\\is off this chart:\\the tolerance\\binds 4$\times$ sooner};

\draw[<-, thick, warnrule] (axis cs:0.50,0.198) -- (axis cs:0.62,0.13)
  node[anchor=west, font=\scriptsize, text=warnrule] {0.202 vs 0.200};

\end{axis}
\end{tikzpicture}

\caption{Three measured points and one straight line. Lag is proportional to
commanded velocity, so the path tolerance is a speed limit in disguise.}
\label{fig:lag}
\end{figure}

The three points fall on a line: $\mathrm{lag} \approx 0.4\,\mathrm{s} \times v$.
The arm runs a roughly constant \emph{time} behind its setpoint, which is what a
series-elastic arm under a proportional controller does.

That reframes the tolerance completely. \texttt{path\_tolerance\_rad} $= 0.2$
divided by 0.4\,s gives about 0.5\,rad/s, and that is the practical speed ceiling
of this workspace on this robot. The shim's 2.0\,rad/s per-cycle clamp is
unreachable --- the tolerance binds four times sooner. Anybody reading the clamp
as the speed limit is reading the wrong number.

\begin{notebox}[Two things this does \emph{not} license]
Do not tighten \texttt{path\_tolerance\_rad}: it is the speed limit, not a safety
margin with slack in it. And do not raise it at a desk. Whether 0.3--0.4 is safe
is an open experiment for a supervised session, not a default someone changes in
a config file.
\end{notebox}

A related result from the same session: \texttt{speed\_ratio} had no measurable
effect across 0.1, 0.2 and 0.3. The client's own interpolation is the binding
constraint, so \texttt{duration} is the knob that changes speed.

\section{MoveIt on hardware}
\label{sec:moveit-hardware}

MoveIt planned and executed on the robot on 2026-07-24. The 2026-07-25 re-run,
with eleven goals per arm, went differently:

\begin{table}[h]
\centering
\small
\begin{tabular}{@{}lrrl@{}}
\toprule
\textbf{Arm} & \textbf{Completed} & \textbf{Aborted} & \textbf{Failure} \\
\midrule
Left  & 3 of 11 & 8 & \texttt{PATH\_TOLERANCE\_VIOLATED} on \texttt{left\_w0} / \texttt{left\_e0} \\
Right & 9 of 11 & 2 & --- \\
\bottomrule
\end{tabular}
\end{table}

Both arms were driven by the same \texttt{both\_arms} plans, and every left-arm
abort landed at exactly the tolerance limit on a wrist or elbow joint. The
mechanism is consistent with Section~\ref{sec:speed}: MoveIt plans at default
velocity scaling, the wrist joints permit 4.0\,rad/s against the shoulder's 1.5,
so a time-parameterised plan drives the wrists fastest and they are the first to
exceed a lag budget that corresponds to about 0.5\,rad/s.

\subsection{Workaround}

On this robot, drive MoveIt at \texttt{Velocity Scaling: 0.1} in the
MotionPlanning panel, and plan single-arm groups rather than \texttt{both\_arms}.

\subsection{What is not yet known}

Whether this is an asymmetric \emph{plan} or an asymmetric \emph{arm} is
undecided. A \texttt{both\_arms} plan is not required to distribute motion
symmetrically, so the left arm may simply have drawn the faster half. The
decisive experiment is to command the same joint on each arm --- the
\texttt{joint} parameter of the direct client exists for this --- at matched
speeds, with the recorder running. The recorder was not running during the
MoveIt re-run, which is precisely why the question is still open.

\section{Joint states arrive from two publishers}
\label{sec:jointstate-merge}

\texttt{/robot/joint\_states} is published by two nodes: arm joints from one,
gripper joints from another. A client keeping only the most recent message sees
no arm joints in roughly a third of messages.

The symptom is not a clean error. It looks like intermittent missing data,
timeouts that come and go, and readings that are stale for no visible reason ---
the sort of thing that gets blamed on the network or the robot. Clients on the
hardware path merge positions by name across messages and tolerate an incomplete
view rather than treating it as a fault.

There is a related defect in the bridge worth knowing about before adding topics:
its subscriber negotiates with only the \emph{first} publisher the master lists,
so on a multi-publisher topic it silently receives a subset. Measured effect:
\texttt{/robot/joint\_states} arrives at about 100\,Hz instead of its real
138\,Hz, and gripper joint states never arrive at all.
